% ============================================================================
%  Capitolo 1 -- PRG32 e PRG32-QT: allestire l'officina          (Lezione 01)
% ============================================================================
\chapter{PRG32 e PRG32-QT: allestire l'officina}
\label{ch:setup}

\epigraphline{Datemi sei ore per abbattere un albero, e passerò le prime
quattro ad affilare l'ascia.}{attribuita ad Abraham Lincoln}

\begin{chapgoals}
Alla fine di questo capitolo avrete
\begin{itemize}
\item un quadro chiaro dei pezzi di \prg: runtime, cartuccia, strumenti, host, Store;
\item \prgqt installato sul vostro portatile, per eseguire qualunque cartuccia e guardarci dentro;
\item la toolchain \riscv installata e verificata, per compilare cartucce vostre;
\item una cartuccia compilata da voi, in esecuzione in \prgqt e fermata nel debugger;
\item un elenco di che cosa fare quando qualcosa non funziona.
\end{itemize}
\end{chapgoals}

Il primo lavoro di un artigiano in una nuova officina è mettere gli attrezzi
dove la mano se li aspetta. Questo capitolo è quel lavoro. Non contiene teoria,
ed è l'unico capitolo del libro il cui contenuto invecchierà: i numeri di
versione sono quelli di ottobre 2026, e prima di seguire i passi conviene
controllare le pagine dei rilasci dei progetti. Tutto il resto del libro
dipende dall'ultima riga del riepilogo qui sopra: il ciclo che consiste nel
compilare una cartuccia, inviarla a un host e fermarla nel debugger si ripete
in ogni capitolo che segue.

I comandi della Sezione~\ref{sec:setup:build} sono stati eseguiti per questo
libro su macOS con \prgqt. Le installazioni per Windows e Linux seguono le
istruzioni ufficiali dei progetti e non sono state eseguite per questo testo.

% ----------------------------------------------------------------------------
\section{Quattro livelli: installate ciò che vi serve}

Non tutto va installato, e non tutto subito. La
Tabella~\ref{tab:setup:levels} ordina i pezzi secondo la necessità. Il corso, e
il libro, si possono seguire con i soli livelli 1 e 2: \emph{non serve alcuna
scheda}.

\begin{table}[htb]
\centering\small
\caption{Quattro livelli di installazione.}\label{tab:setup:levels}
\begin{tabularx}{\textwidth}{@{}cP{0.24\textwidth}LP{0.22\textwidth}@{}}
\toprule
\textbf{Livello} & \textbf{Installate} & \textbf{Potete} & \textbf{Serve}\\
\midrule
1 & \prgqt & eseguire cartucce, scaricarle dallo Store, usare il debugger & a tutti, subito\\
2 & ESP-IDF e il repository di \prg & compilare cartucce vostre da assembly o C & a tutti, prima del Capitolo~\ref{ch:firstcontact}\\
3 & l'emulatore QEMU & eseguire il firmware vero sul portatile & facoltativo\\
4 & il kit: una scheda ESP32-C6 & eseguire sull'hardware vero & facoltativo\\
\bottomrule
\end{tabularx}
\end{table}

Il livello~1 richiede cinque minuti. Il livello~2 è un download voluminoso e
richiede da mezz'ora a un'ora: conviene farlo a casa, con una buona
connessione. I livelli 3 e 4 sono per i curiosi, e per il laboratorio.

% ----------------------------------------------------------------------------
\section{La piattaforma, in dettaglio}

\subsection{Un runtime, un formato di file, un insieme di strumenti}

Il nome \prg copre quattro cose.

\begin{itemize}
\item Un \textbf{firmware residente}, detto anche \term{runtime}: il programma
  che possiede lo schermo, i pulsanti e il suono.
\item Un \textbf{formato di cartuccia}: un gioco come un unico piccolo file,
  \fname{.prg32}.
\item Uno \textbf{strumento} sul portatile: il comando \code{python3 -m prg32},
  che compila, carica e confeziona.
\item Un \textbf{Cartridge Store}, dove si condividono i giochi finiti.
\end{itemize}

Il codice di una cartuccia gira come vero codice macchina \riscv a 32 bit,
RV32IMAC, e il contratto tra un gioco e il runtime sono tre funzioni:
\code{init}, \code{update} e \code{draw}. Il runtime chiama \code{init} una
volta e poi \code{update} e \code{draw} a ogni frame; il
Capitolo~\ref{ch:firstcontact} parte da questo contratto.

\subsection{Da un file sorgente allo schermo}

La Figura~\ref{fig:setup:flow} segue un programma dall'editor allo schermo. Il
sorgente, in assembly o in C, viene tradotto in codice macchina dalla toolchain
\riscv; il compilatore è \code{riscv32-esp-elf-gcc}, installato da ESP-IDF, il
kit di sviluppo di Espressif. Il risultato viene impacchettato in una
cartuccia, fatta di codice, dati e un'intestazione. La cartuccia viene inviata
attraverso la rete a un \term{host}, dove il runtime la carica e comincia a
chiamare le sue tre funzioni.

\begin{figure}[htb]
\centering
\begin{tikzpicture}[node distance=5mm]
  \node[blkacc,text width=21mm,minimum height=15mm] (a) {\texttt{verm1.S}\\[1pt]\footnotesize il vostro sorgente};
  \node[blkamb,text width=21mm,minimum height=15mm,right=of a] (b) {compilatore\\[1pt]\footnotesize toolchain \riscv};
  \node[blkvio,text width=23mm,minimum height=15mm,right=of b] (c) {\texttt{verm1.prg32}\\[1pt]\footnotesize una cartuccia};
  \node[blkgrn,text width=19mm,minimum height=15mm,right=of c] (d) {caricamento\\[1pt]\footnotesize su un host};
  \node[blkred,text width=21mm,minimum height=15mm,right=of d] (e) {runtime\\[1pt]\footnotesize chiama \texttt{init}, \texttt{update}, \texttt{draw}};
  \foreach \x/\y in {a/b,b/c,c/d,d/e} \draw[->,thick] (\x) -- (\y);
  \draw[decorate,decoration={brace,mirror,amplitude=4pt}] ([yshift=-2mm]a.south west) -- node[below=5pt,lbl]{sul vostro portatile: serve il livello 2} ([yshift=-2mm]c.south east);
  \draw[decorate,decoration={brace,mirror,amplitude=4pt}] ([yshift=-2mm]d.south west) -- node[below=5pt,lbl]{un host: \prgqt, QEMU o la scheda} ([yshift=-2mm]e.south east);
\end{tikzpicture}
\caption{Da un file sorgente allo schermo.}\label{fig:setup:flow}
\end{figure}

Una cartuccia viene caricata all'indirizzo \hex{40800000} e raggiunge il
runtime soltanto attraverso una tabella di funzioni, l'\term{application binary
interface} (ABI). È la ragione per cui lo stesso file gira su tre host
diversi: non tocca mai direttamente l'hardware, e ogni host fornisce la stessa
tabella.

\subsection{Il kit: una console su breadboard}

L'hardware di \prg (Figura~\ref{fig:setup:kit}) è costruito attorno a un
ESP32-C6 di Espressif: un core \riscv a 160~MHz, con Wi-Fi. Un display a colori
da 2,8~pollici di 320 per 240 pixel (un controller ILI9341 su un bus SPI)
mostra i giochi, che disegnano in un viewport di 320 per 200. Uno o due piccoli
amplificatori audio (MAX98357A, su un bus I2S), un joystick e dei pulsanti sono
collegati direttamente ai piedini del chip. Un gioco raggiunge la scheda via
Wi-Fi: la scheda crea una rete chiamata \code{PRG32}, e la cartuccia vi viene
caricata. L'elenco dei componenti e i collegamenti sono nella cartella
\fname{docs/hardware} del repository.

\begin{figure}[htb]
\centering
\includegraphics[width=0.7\textwidth]{kit_breadboard.png}
\caption{Il kit \prg su breadboard.}\label{fig:setup:kit}
\end{figure}

\subsection{Tre modi di eseguire la stessa cartuccia}

Il file della cartuccia è lo stesso su ogni host; cambia il modo in cui le sue
istruzioni vengono eseguite e ciò che si può osservare
(Tabella~\ref{tab:setup:hosts}).

\begin{table}[htb]
\centering\small
\caption{Tre host per la stessa cartuccia.}\label{tab:setup:hosts}
\begin{tabularx}{\textwidth}{@{}P{0.19\textwidth}LLL@{}}
\toprule
 & \textbf{Scheda ESP32-C6} & \textbf{QEMU di Espressif} & \textbf{\prgqt}\\
\midrule
Come gira il codice & in modo nativo sul chip & in modo nativo su un ESP32-C3 emulato & guest RV32IMAC interpretato\\
Serve una toolchain & sì (ESP-IDF) & sì (ESP-IDF) & no, per giocare e fare debug\\
La cartuccia arriva con & caricamento via Wi-Fi, Store & scrittura nell'immagine della flash & importazione di file, Store, caricamento HTTP\\
Esecuzione passo-passo & con JTAG e GDB & con GDB & debugger integrato\\
Gira su & il kit hardware & Windows, Linux, macOS & desktop, Raspberry Pi, iOS, Android, TV\\
\bottomrule
\end{tabularx}
\end{table}

La regola pratica è: usate la scheda quando volete la verità sul \emph{tempo},
e \prgqt quando volete la verità su un \emph{registro}. L'emulatore modella un
ESP32-C3 e non un C6 perché Espressif mantiene per quel chip un emulatore
maturo con grafica; esegue le stesse istruzioni \riscv, e la cartuccia non se
ne accorge.

\subsection{PRG32-QT: la console sul portatile}

\prgqt è un'applicazione per Windows, macOS, Linux, Raspberry Pi, telefoni,
tablet e televisori. Contiene un interprete di RV32IMAC, che esegue una
cartuccia istruzione per istruzione e le offre le stesse funzioni del firmware.
Ha tre velocità, scelte in \emph{Settings}: \emph{Accurate}, temporizzata come
l'ESP32-C6, che è quella predefinita; \emph{Optimal}, a 30 frame al secondo; e
\emph{Unlimited}, veloce quanto lo è il portatile. Il profilo accurato
attribuisce a ogni istruzione e a ogni chiamata un costo in cicli di un clock a
160~MHz: non è un modello esatto al ciclo del chip, ma è abbastanza vicino da
far vedere quando un frame è in ritardo. \prgqt contiene anche un browser dello
Store e il debugger descritto nel Capitolo~\ref{ch:intro}.

I comandi della console corrispondono alla tastiera come nella
Tabella~\ref{tab:setup:keys}. Un controller di gioco USB o Bluetooth viene
riconosciuto da \prgqt quando lo si collega.

\begin{table}[htb]
\centering\small
\caption{I comandi, su una tastiera.}\label{tab:setup:keys}
\begin{tabular}{@{}lll@{}}
\toprule
\textbf{Sul kit} & \textbf{In \prgqt} & \textbf{In QEMU}\\
\midrule
Joystick & tasti freccia, oppure \key{W A S D} & gli stessi, battuti nel terminale\\
Pulsante A & \key{Z} o \key{J} & \key{Z} o \key{J}\\
Pulsante B & \key{X} o \key{K} & \key{X} o \key{K}\\
Start / Select & \key{Invio} o \key{Spazio} & \key{Invio} o \key{Spazio}\\
Uscire dal gioco & \key{Esc}: ritorno a Setup & chiudere QEMU\\
\bottomrule
\end{tabular}
\end{table}

\begin{prgwarn}
QEMU legge la tastiera dal terminale che lo ha lanciato, non dalla finestra con
l'immagine. È la classica sorpresa della prima sessione.
\end{prgwarn}

% ----------------------------------------------------------------------------
\section{Passo 1: installare PRG32-QT}

Cinque minuti; nessun compilatore, nessun account. Dalla pagina
\url{https://github.com/riscv-prg32/PRG32-QT}, sotto \emph{Releases}, prendete
l'ultima versione del file per il vostro sistema
(Tabella~\ref{tab:setup:files}). Un file \fname{SHA256SUMS-platforms.txt} sulla
stessa pagina elenca le somme di controllo; lì ci sono anche le versioni per
Android e per i televisori.

\begin{table}[htb]
\centering\small
\caption{Gli installer di \prgqt (versione 0.4.1, ottobre 2026; i nomi seguono
questo schema con il numero corrente).}\label{tab:setup:files}
\begin{tabular}{@{}lll@{}}
\toprule
\textbf{Sistema} & \textbf{File} & \textbf{Tipo}\\
\midrule
Windows 10 o 11 & \fname{PRG32-QT-0.4.1-Windows-X64-Setup.exe} & installer\\
macOS, Apple Silicon & \fname{PRG32-QT-0.4.1-macOS-ARM64.dmg} & immagine disco\\
macOS, Intel & \fname{PRG32-QT-0.4.1-macOS-Intel-X64.dmg} & immagine disco\\
Linux, PC a 64 bit & \fname{PRG32-QT-0.4.1-Linux-X64.deb} & pacchetto Debian\\
Raspberry Pi OS, 64 bit & \fname{PRG32-QT-0.4.1-Raspberry-Pi-OS-ARM64.deb} & pacchetto Debian\\
\bottomrule
\end{tabular}
\end{table}

\paragraph{Windows.} Scaricate il file che termina con
\fname{Windows-X64-Setup.exe} ed eseguitelo. Se Windows mostra un avviso blu di
SmartScreen, scegliete \emph{Ulteriori informazioni} e poi \emph{Esegui
comunque}. Accettate la cartella proposta: l'installer porta con sé le librerie
Qt e il runtime del compilatore, e non c'è altro da installare. Avviate
\emph{PRG32} dal menu Start; se il firewall lo chiede, autorizzatelo sulle reti
private. La domanda compare perché \prgqt ascolta sulla porta 8080 le cartucce
inviate dagli strumenti.

\paragraph{macOS.} Dal menu Apple scegliete \emph{Informazioni su questo Mac}:
«Chip Apple M\ldots» significa ARM64, «Processore Intel» significa X64. Aprite
il \fname{.dmg} e trascinate \emph{PRG32} sulla cartella Applicazioni. Al primo
avvio, se macOS dice che lo sviluppatore non può essere verificato
(l'applicazione è distribuita fuori dall'App Store), fate clic con il tasto
destro su \emph{PRG32} e scegliete \emph{Apri}, oppure andate in
\emph{Impostazioni di Sistema, Privacy e sicurezza, Apri comunque}. Accettate
le connessioni in entrata e l'accesso alla rete locale quando viene chiesto:
servono all'interfaccia di dispositivo sulla porta 8080 e al suo annuncio sulla
rete locale.

\paragraph{Linux.} Il pacchetto richiede Qt~6.5 o successivo, che Debian~13 e
Ubuntu~25.04 forniscono; sui sistemi più vecchi (Ubuntu~24.04 ha Qt~6.4) il
programma si compila da sorgente (Listato~\ref{lst:setup:qtlinux}). Su
Raspberry Pi OS usate il pacchetto ARM64, oppure lo script
\fname{scripts/build-raspbian.sh}.

\shfile[caption={\prgqt su Linux.},label=lst:setup:qtlinux]{code/l01/c06.sh}

\begin{prgtry}[Eseguite: la vostra prima cartuccia, dallo Store]
\begin{enumerate}[nosep]
\item Avviate \prgqt: compare la schermata \emph{Setup}, la stessa che il
  firmware mostra sul kit.
\item Scegliete \emph{Browse Store}: viene scaricato il catalogo del Cartridge
  Store.
\item Scegliete un gioco, scaricatelo, e scegliete \emph{Run Cartridge}.
\item Giocate con i tasti freccia, \key{Z} e \key{X}. \key{Esc} vi riporta a
  Setup.
\item Aprite \emph{Settings}: guardate i tre profili di prestazioni e le scelte
  di disposizione.
\end{enumerate}
\textbf{Risultato atteso.} Un gioco dello Store gira sul vostro portatile
(Figura~\ref{fig:setup:screens}). Questo è il livello~1: da ora potete eseguire
qualunque cartuccia, e guardarci dentro. Se l'elenco dello Store resta vuoto,
il portatile non raggiunge Internet oppure c'è di mezzo un proxy:
\emph{Settings} mostra l'indirizzo dello Store.
\end{prgtry}

\begin{figure}[htb]
\centering
\begin{subfigure}[b]{0.32\textwidth}\includegraphics[width=\textwidth]{menu_main.png}\caption{Il menu Setup}\end{subfigure}\hfill
\begin{subfigure}[b]{0.32\textwidth}\includegraphics[width=\textwidth]{prg32qt_store.png}\caption{Browse Store}\end{subfigure}\hfill
\begin{subfigure}[b]{0.32\textwidth}\includegraphics[width=\textwidth]{prg32qt_settings.png}\caption{Settings}\end{subfigure}
\caption{Che cosa dovreste vedere: le stesse tre schermate su ogni sistema.}
\label{fig:setup:screens}
\end{figure}

% ----------------------------------------------------------------------------
\section{Passo 2: installare la toolchain}

\subsection{Un'idea dietro ogni installazione}

Il compilatore per la famiglia ESP32 arriva con ESP-IDF, il kit di sviluppo di
Espressif \citep{espidf}. \prg richiede ESP-IDF~5.4 o successivo, installato
per due chip: \code{esp32c6}, il kit, ed \code{esp32c3}, l'emulatore. Dopo di
che, un solo comando fa tutto: \code{python3 -m prg32}.

\begin{prgwarn}[ESP-IDF non è attivo da solo]
Ogni nuovo terminale deve accendere ESP-IDF prima che gli strumenti funzionino.
Quasi ogni problema di installazione è ESP-IDF non attivo \emph{nel terminale
in cui state scrivendo}. Il sintomo è sempre lo stesso messaggio:
\begin{center}\code{idf.py: command not found}\end{center}
(oppure \code{riscv32-esp-elf-gcc} non trovato). Su Windows, attivare significa
aprire la \emph{ESP-IDF PowerShell}; su macOS e Linux, eseguire
\fname{export.sh}.
\end{prgwarn}

\subsection{Windows}

Installate Git for Windows e Python~3, poi eseguite l'ESP-IDF Tools Installer e
lavorate nella finestra di PowerShell che esso crea
(Listato~\ref{lst:setup:win}). Usate la ESP-IDF PowerShell, mai un terminale
qualunque; su Windows il comando è \code{python}, non \code{python3};
\code{pip} installa nell'ambiente di ESP-IDF, ed è voluto. Evitate le cartelle
con spazi o lettere accentate nel percorso. Chi preferisce un editor può usare
Visual Studio Code con l'estensione Espressif ESP-IDF, che installa la stessa
toolchain.

\shfile[caption={La toolchain su Windows.},label=lst:setup:win]{code/l01/c09.sh}

\subsection{macOS}

Prima viene Homebrew, perché fornisce \code{git}, \code{cmake} e \code{python}
(Listato~\ref{lst:setup:mac}); \code{git} può chiedere, la prima volta, di
installare gli strumenti a riga di comando di Xcode. Il clone ricorsivo e
\fname{install.sh} scaricano qualche gigabyte: abbiate pazienza. Il punto prima
di \fname{export.sh} conta, perché fa sì che lo script modifichi \emph{questo}
terminale. La procedura funziona su Apple Silicon e su Intel.

\shfile[caption={La toolchain su macOS.},label=lst:setup:mac]{code/l01/c10.sh}

\subsection{Linux (Debian, Ubuntu)}

I prerequisiti sono una lunga riga di \code{apt}, da copiare per intero
(Listato~\ref{lst:setup:linux}); i passi 2 e 3 sono gli stessi di macOS. Altre
distribuzioni hanno gli stessi pacchetti con altri nomi. Solo per il kit,
l'utente deve appartenere al gruppo \code{dialout}, che dà accesso alla porta
seriale della scheda: \code{sudo usermod -aG dialout \$USER}, poi si esce e si
rientra nella sessione.

\shfile[caption={La toolchain su Linux.},label=lst:setup:linux]{code/l01/c11.sh}

\subsection{A ogni nuovo terminale: accendere ESP-IDF}

\fname{export.sh} agisce sul terminale in cui viene eseguito, e solo su quello.
Una nuova finestra è un nuovo terminale, e va attivata di nuovo. Due alias
risparmiano battute (Listato~\ref{lst:setup:activate}); sono una comodità, non
un obbligo. I comandi \code{prg32} si eseguono dalla cartella di \prg.

\shfile[caption={Attivare ESP-IDF in un nuovo terminale.},label=lst:setup:activate]{code/l01/c12.sh}

\begin{prgtry}[Vedete l'errore una volta, di proposito]
Chiudete il terminale, apritene uno nuovo ed eseguite \code{python3 -m prg32
doctor} \emph{senza} attivare ESP-IDF. Leggete il messaggio. È meglio
incontrarlo qui che nel mezzo di un tutorial.
\end{prgtry}

\subsection{Il doctor}

Il comando \code{doctor} controlla l'installazione
(Listato~\ref{lst:setup:doctor}). La riga su \code{idf.py} dice che ESP-IDF è
attivo; quella su \code{riscv32-esp-elf-gcc} che il compilatore \riscv è
raggiungibile; quelle sulle partizioni che la cartella di \prg è completa. I
quattro slot per le cartucce, da \code{cart0} a \code{cart3}, sono regioni
della memoria flash della scheda, di 131\,072 byte ciascuna. Ogni riga marcata
\code{ERROR} dice che cosa fare. Se il doctor passa, il livello~2 è fatto.

\txtfile[caption={Il doctor (messaggi di ottobre 2026).},label=lst:setup:doctor]{code/l01/c13.sh}

% ----------------------------------------------------------------------------
\section{Passo 3: la vostra prima compilazione}\label{sec:setup:build}

\subsection{Compilare una cartuccia dagli esempi}

Il repository contiene giochi completi. Compilarne uno dimostra che il
compilatore funziona prima ancora di aver scritto una sola riga
(Listato~\ref{lst:setup:build}).

\txtfile[caption={Compilare Pong, in C, dagli esempi.},label=lst:setup:build]{code/l01/c14.sh}

Il comando \code{cartridge build} prende un sorgente e produce un file
\fname{.prg32}. L'opzione \code{-{}-entry-prefix} dà il nome che nel sorgente
precede \code{\_init}, \code{\_update} e \code{\_draw}; \code{-{}-portable}
chiede una cartuccia che giri su ogni host, e non su un solo firmware. Un Pong
intero sono tremila byte. Nel Capitolo~\ref{ch:firstcontact} il sorgente sarà
quello del lettore, in assembly, e il comando sarà lo stesso. Su Windows
scrivete \code{python} e mettete il comando su una sola riga.

\subsection{Inviarla a PRG32-QT}

\prgqt risponde sulla porta 8080 esattamente come una scheda
(Listato~\ref{lst:setup:send}). Il comando \code{runtime} stampa lo stato
dell'host ed è una prova rapida della connessione; \code{upload-and-run}
memorizza la cartuccia in uno slot e la avvia. Il comando dice \code{esp32c6}
perché, per lo strumento, \prgqt \emph{è} una scheda. Se \code{localhost} non
risponde, leggete l'indirizzo mostrato nella schermata Setup di \prgqt e usate
quello. Senza gli strumenti, lo stesso risultato si ottiene in \prgqt
importando il file \fname{pong.prg32} e scegliendo \emph{Run Cartridge}; con
\code{-{}-url http://192.168.4.1} lo stesso comando invia il gioco a una scheda
vera.

\txtfile[caption={Inviare la cartuccia a \prgqt, in esecuzione sullo stesso calcolatore.},label=lst:setup:send]{code/l01/c15.sh}

\subsection{Guardare dentro: il debugger}

\begin{prgtry}[Aprite la macchina]
\begin{enumerate}[nosep]
\item In \prgqt aprite \emph{Settings} e attivate \emph{Run cartridges in
  RISC-V debug mode}.
\item Eseguite di nuovo la cartuccia: il gioco è a sinistra, la macchina a
  destra (Figura~\ref{fig:setup:debug}).
\item Premete \key{Pause} (il pulsante con due barre): il gioco si blocca, e
  nel listato è evidenziata un'istruzione.
\item Premete \key{Step} (il triangolo con una barra) qualche volta:
  l'evidenziazione si sposta, e alcuni registri cambiano valore.
\item Premete \key{Update}, poi \key{Draw}: il listato salta a quelle due
  funzioni della cartuccia.
\item Premete \key{Resume} (il triangolo): il gioco prosegue. Abbassate la
  velocità a 0.1x e guardatelo al rallentatore.
\end{enumerate}
\end{prgtry}

I tre pulsanti di trasporto del debugger sono disegnati come su un lettore
video; questo libro li chiama \key{Pause}, \key{Resume} e \key{Step}.
Questo è solo un primo sguardo: il prossimo capitolo spiega che cosa sono i
registri. Qui importa che il debugger si apra e risponda. La modalità debug
resta attiva finché non viene disattivata.

\begin{figure}[htb]
\centering
\includegraphics[width=0.95\textwidth]{prg32qt_debugger_details.png}
\caption{La macchina, aperta: l'assembly con l'istruzione corrente, i 32
registri, la memoria, e i contatori del frame. Se vedete questo con la vostra
cartuccia, siete pronti per il resto del libro.}
\label{fig:setup:debug}
\end{figure}

% ----------------------------------------------------------------------------
\section{Facoltativo: QEMU e il kit}

\subsection{Livello 3: l'emulatore QEMU}

QEMU emula un chip intero ed esegue il firmware vero
(Listato~\ref{lst:setup:qemu}). Il comando \code{qemu build} compila il
firmware e la prima volta richiede parecchi minuti, perché scarica dei
componenti; \code{qemu upload} scrive una cartuccia nell'immagine della flash,
e mentre gira QEMU deve essere chiuso. Se il caricamento fallisce, eseguite
QEMU una volta perché l'immagine della flash esista, chiudetelo, e caricate di
nuovo. Da ottobre 2026 il firmware parte per impostazione predefinita dal menu
\emph{Setup}, su QEMU e sulla scheda: scegliete \emph{Run Cartridge}, poi lo
slot \code{cart0}. Questa procedura viene dalla documentazione di \prg
(\fname{docs/usage/qemu.md}) e non è stata eseguita per questo libro.

\shfile[caption={Livello 3: l'emulatore QEMU.},label=lst:setup:qemu]{code/l01/c18.sh}

\subsection{Livello 4: il kit}

Il firmware viene scritto sulla scheda una volta, attraverso un cavo USB; dopo
di che i giochi arrivano via Wi-Fi e il cavo non serve più
(Listato~\ref{lst:setup:kit}). Il comando di caricamento è lo stesso usato per
\prgqt, con un altro indirizzo. Se la porta seriale non viene trovata la si
indica esplicitamente: \code{COM}$x$ su Windows, \fname{/dev/ttyACM0} su Linux,
\fname{/dev/cu.usbmodem\ldots} su macOS. Uno schermo nero dopo la scrittura
significa di solito che è stato scritto il firmware per QEMU invece di quello
per la scheda. Anche questa procedura viene dalla documentazione e non è stata
eseguita per questo libro.

\shfile[caption={Livello 4: il kit.},label=lst:setup:kit]{code/l01/c19.sh}

% ----------------------------------------------------------------------------
\section{Quando qualcosa non funziona}

La Tabella~\ref{tab:setup:trouble} elenca i sintomi che si incontrano più
spesso; la prima riga copre la maggior parte dei casi. Il file
\fname{docs/usage/troubleshooting.md} del repository contiene l'elenco
completo.

\begin{table}[htb]
\centering\small
\caption{Risoluzione dei problemi.}\label{tab:setup:trouble}
\begin{tabularx}{\textwidth}{@{}P{0.27\textwidth}LL@{}}
\toprule
\textbf{Vedete} & \textbf{Significa} & \textbf{Fate così}\\
\midrule
\code{idf.py: command not found} & ESP-IDF non è attivo in questo terminale & eseguite \fname{export.sh}, o aprite la ESP-IDF PowerShell\\
manca \code{riscv32-esp-elf-gcc} & i compilatori non sono stati installati, o non sono attivi & \code{./install.sh esp32c3,esp32c6}, poi \fname{export.sh}\\
\code{No module named prg32} & non siete nella cartella di \prg & \code{cd} nella cartella che avete clonato\\
upload: connection refused & \prgqt non è in esecuzione, o un firewall blocca la porta 8080 & avviatelo; autorizzatelo; usate l'indirizzo della sua schermata Setup\\
macOS non apre PRG32 & l'applicazione viene da fuori dell'App Store & clic destro, Apri; oppure Privacy e sicurezza, Apri comunque\\
Linux: PRG32 non parte & il Qt del sistema è più vecchio della 6.5 & una distribuzione più recente, o compilate da sorgente\\
QEMU ignora i tasti & la tastiera va al terminale & fate clic sul terminale che ha lanciato QEMU\\
\bottomrule
\end{tabularx}
\end{table}

\begin{prgconcept}[Come chiedere aiuto]
Incollate il comando e il messaggio \emph{intero} come testo. Non mandate mai
una fotografia dello schermo.
\end{prgconcept}

\subsection{Pronti per il prossimo capitolo?}

Cinque controlli, e l'officina è pronta.

\begin{enumerate}
\item \prgqt parte ed esegue un gioco dello Store (passo~1).
\item In un nuovo terminale, dopo aver attivato ESP-IDF, \code{doctor} mostra
  solo \code{OK} (passo~2).
\item \code{cartridge build} produce \fname{pong.prg32} (passo~3).
\item \code{upload-and-run} lo avvia in \prgqt (passo~3).
\item In modalità debug, \key{Pause} e \key{Step} funzionano (passo~3).
\end{enumerate}

% ----------------------------------------------------------------------------
% Generated by tools/textbook/export_listings.js from decks/plan.js and code/l01.S: do not edit.
\section{Il programma di questo capitolo}\label{sec:lab:l01}

I listati di questo capitolo sono parti di un solo programma,
\fname{code/l01.S} nel repository del corso, stampato per intero nel
Listato~\ref{lst:lab:l01}. I test automatici del materiale del corso lo
compilano e lo eseguono sull'esecutore senza interfaccia di \prgqt{}, con
ingressi programmati, e controllano il frame; la
Figura~\ref{fig:lab:l01} è il frame scritto dal test.

\begin{description}[style=nextline,leftmargin=1.5em]
\item[Sullo schermo.] Quattro barre verticali, rossa, verde, blu e bianca, e le parole PRG32 OK.
\item[Che cosa premere.] Niente: il programma va da solo.
\end{description}

\begin{figure}[htb]
\centering
\includegraphics[width=0.5\textwidth]{lab_l01.png}
\caption{\fname{l01.S} su \prgqt{}.}\label{fig:lab:l01}
\end{figure}

\asmfile[caption={Il programma completo del capitolo, \fname{code/l01.S}.},label=lst:lab:l01]{code/l01/program.S}

Per eseguirlo su \prgqt{} i comandi sono quelli del
Listato~\ref{lst:lab:l01:run}; il Paragrafo~\ref{sec:prg:testing} dà gli
stessi passi per QEMU e per la scheda.

\begin{shell}[caption={Compilare \fname{l01.S} ed eseguirlo su \prgqt{}.},label=lst:lab:l01:run]
# 1. start PRG32-QT on this computer (it listens on port 8080)

# 2. in the PRG32 folder, with ESP-IDF active: build the cartridge
python3 -m prg32 cartridge build ~/arc12/code/l01.S --portable \
    --entry-prefix l01 --name l01 --out l01.prg32

# 3. send it to PRG32-QT and start it
python3 -m prg32 esp32c6 upload-and-run l01.prg32 \
    --url http://localhost:8080 --slot cart0

# 4. look, press, compare with the previous slide
\end{shell}


% ----------------------------------------------------------------------------
\section{Verificate di aver capito}

\begin{quiz}{In un nuovo terminale, \code{python3 -m prg32 doctor} dice:
\code{idf.py missing}. Qual è la causa più probabile?}
\item ESP-IDF va installato di nuovo
\item ESP-IDF non è attivo in questo terminale
\item \prgqt non è in esecuzione
\item la scheda non è collegata
\end{quiz}

\begin{quiz}{Quale host permette di eseguire una cartuccia e farne il debug senza installare un compilatore?}
\item la scheda ESP32-C6
\item QEMU
\item \prgqt
\item nessuno
\end{quiz}

\begin{quiz}{Perché lo stesso \fname{pong.prg32} gira sulla scheda, su QEMU e su \prgqt?}
\item viene ricompilato a ogni avvio
\item contiene tre programmi
\item chiama il runtime attraverso la stessa tabella di funzioni su ogni host
\item \prgqt lo traduce in C
\end{quiz}

\section{Esercizi}

\begin{exercise}[l'officina sopravvive a un riavvio]
\begin{enumerate}[nosep]
\item Completate i cinque controlli qui sopra sul vostro portatile.
\item Scaricate tre cartucce dallo Store, e in ciascuna trovate con il debugger
  la prima istruzione di \code{update}.
\item Compilate uno degli altri esempi in \fname{examples/games}, ed
  eseguitelo.
\item Chiudete tutto, riavviate il calcolatore, e tornate a una cartuccia in
  esecuzione senza questo libro.
\item Facoltativo: installate QEMU ed eseguite lì la stessa cartuccia.
\end{enumerate}
\expected{La vostra officina sopravvive a un riavvio, e sapete aprirla senza
aiuto. Il quarto compito è la vera prova: mostra se l'attivazione di ESP-IDF è
stata capita.}
\end{exercise}

\section{Glossario}

\begin{glossarytable}
toolchain & catena di strumenti di sviluppo\\
cartridge & cartuccia\\
emulator & emulatore\\
to build & compilare (costruire)\\
host & macchina ospite\\
firmware & firmware\\
runtime & ambiente di esecuzione\\
interpreter & interprete\\
to flash & scrivere nella memoria flash\\
breakpoint & punto di interruzione\\
\end{glossarytable}

\section{Riepilogo e letture}

\prg è un runtime, un formato di cartuccia, uno strumento e uno Store; una
cartuccia gira senza modifiche sulla scheda, su QEMU e su \prgqt. \prgqt basta
per giocare e per fare debug; ESP-IDF e il repository di \prg servono per
compilare. ESP-IDF va attivato in ogni nuovo terminale: ricordatelo e la
maggior parte dei problemi scompare. Avete compilato una cartuccia, l'avete
inviata a \prgqt e l'avete fermata nel debugger: è il ciclo di tutto il libro.

\paragraph{Letture.} Il \fname{README} e la cartella \fname{docs/usage} del
repository di \prg \citep{prg32repo}; il \fname{README} e la documentazione di
\prgqt \citep{prg32qt}; i capitoli sull'installazione e su \prgqt di
\emph{Playing with RISC-V} \citep{prg32textbook}. L'Appendice~\ref{app:prg32}
di questo libro è un riferimento della piattaforma e dei suoi comandi.